
\subsection{切线放缩}

\begin{theorem}[切线放缩]
    \begin{tikzpicture}
        \begin{axis}[
                width=7cm, height=7cm,
                title={\textbf{函数 $y=e^x$, $y=\ln x$ 及其相关切线图}},
                xlabel={$x$},
                ylabel={$y$},
                axis lines=middle,
                grid=major,
                xmin=-3.5, xmax=4.5,
                ymin=-3.5, ymax=6.5,
                legend pos=outer north east,    % 图例显示在图右侧
                legend cell align={left},
                enlarge y limits={upper, value=0.1} % Y轴上方留出一些空间
            ]

            % --- 1. 绘制核心函数 y=e^x 和 y=ln(x) ---
            % 使用较粗的线条突出显示
            \addplot[
                domain=-3.5:2, samples=150, % 定义域和采样点数
                color=red, very thick       % 颜色和线宽
            ] {exp(x)};
            \addlegendentry{$y=e^x$}

            \addplot[
                domain=0.02:4.5, samples=150,
                color=blue, very thick
            ] {ln(x)};
            \addlegendentry{$y=\ln x$}

            % --- 2. 绘制对称轴 y=x ---
            % 使用虚线表示参考线
            \addplot[
                domain=-3.5:4.5,
                color=black, dashed
            ] {x};
            \addlegendentry{$y=x$ (对称轴)}

            % --- 3. 绘制 y=e^x 的切线 ---
            % 使用较细的线条
            \addplot[
                domain=-3:3,
                color=orange, thick
            ] {x+1};
            \addlegendentry{$y=x+1$ (在 $x=0$ 处与 $e^x$ 相切)}

            \addplot[
                domain=-0.5:2.2,
                color=purple, thick
            ] {exp(1)*x}; % e*x
            \addlegendentry{$y=ex$ (在 $x=1$ 处与 $e^x$ 相切)}

            % --- 4. 绘制 y=ln(x) 的切线 ---
            % 使用较细的线条
            \addplot[
                domain=-2.5:4.5,
                color=teal, thick
            ] {x-1};
            \addlegendentry{$y=x-1$ (在 $x=1$ 处与 $\ln x$ 相切)}

            \addplot[
                domain=-3.5:4.5,
                color=brown, thick
            ] {x/exp(1)}; % (1/e)*x
            \addlegendentry{$y=\frac{1}{e}x$ (在 $x=e$ 处与 $\ln x$ 相切)}

        \end{axis}
    \end{tikzpicture}
\end{theorem}

\subsubsection{使用切线放缩证明不等式}

\begin{example}
    （2022 山东菏泽高二期末）已知 $a=\frac{9}{10}, b=\mathrm{e}^{-\frac{1}{9}}, c=1+\ln \frac{10}{11}$ ，则 （\quad）
    \begin{tasks}(4)
        \task $a<b<c$
        \task $b<a<c$
        \task $c<b<a$
        \task $c<a<b$
    \end{tasks}
\end{example}

\v{12em}

\begin{example}
    （2022 四川凉山高二期末）已知 $a=\mathrm{e}^{0.01}, b=1.01, c=1-\ln \frac{100}{101}$ ，则 （\quad）
    \begin{tasks}(4)
        \task $c>a>b$
        \task $a>c>b$
        \task $a>b>c$
        \task $b>a>c$
    \end{tasks}
\end{example}

\v{12em}

\begin{example}
    （2022 四川绵阳高二期末）若 $a=\ln \frac{8}{7}, b=\frac{1}{8}, c=\ln \frac{7}{6}$ ，则 （\quad）
    \begin{tasks}(4)
        \task $a<c<b$
        \task $c<a<b$
        \task $c<b<a$
        \task $b<a<c$
    \end{tasks}
\end{example}

\v{12em}

\begin{example}
    （2022 全国甲卷高考 12 题）已知 $a=\frac{31}{32}, b=\cos \frac{1}{4}, c=4 \sin \frac{1}{4}$ ，则 （\quad）
    \begin{tasks}(4)
        \task $c>b>a$
        \task $b>a>c$
        \task $a>b>c$
        \task $a>c>b$
    \end{tasks}
\end{example}

\v{12em}

\begin{example}
    （2022 吉林长春市二中高二期末）已知 $a=\cos \frac{1}{5}, b=\frac{49}{50}, c=5 \sin \frac{1}{5}$ ，则 （\quad）
    \begin{tasks}(4)
        \task $b>a>c$
        \task $c>b>a$
        \task $b>c>a$
        \task $c>a>b$
    \end{tasks}
\end{example}


\subsubsection{使用切线放缩比较特殊值}
